%% kapa: exam=2 round=25 year=2014 period=1 subject=appraisal_practice subject_ko=감정평가실무 seq=3 number=3 points=20
%% source: Q-Net 감정평가사 기출문제 (cst00309) artlSeq=5109209 file=extracted/1008/5109209/2014년 제25회 감정평가사 2차시험 기출문제/1교시(감정평가실무).hwp
%% conversion: kordoc 4.24.51 -> tools/gen2.py (auto); review=proofread-sonnet; flags=tikz
\begin{longquestion}{【문제 3】}
A획지 소유자는 B획지 소유자의 토지 및 전자공장 1동을 인수하여 자신의 토지와 합병한 후, 재가동하고자 한다. 이에 감정평가사 甲에게 토지와 건물의 적정매입 금액 산정을 의뢰하였다. 주어진 자료에 의거 B획지에 소재한 토지와 공장의 가격을 다음 물음에 의거하여 산정하시오. (20점) \par
\medskip\noindent 물음 1) A와 B토지의 합병으로 인한 증분가치를 구하고, B토지의 적정매입가격을 기여도비율에 의한 차액배분법의 논리로 산정하시오. (합병에 따른 제반 비용은 고려치 아니하기로 함) (10점) \par
\medskip\noindent 물음 2) B토지 지상의 공장에 대한 적정한 가격을 원가법에 의한 적산가격으로 산정하되, 재조달원가는 소유자 제시자료에 의한 직접법에 따라, 감가수정은 설비의 특성을 고려한 정률법에 의하시오. (10점) \par
\textbf{<자료 1> 토지의 상황} \par
각 획지의 형태는 다음의 그림과 같다. 각각의 토지가치를 파악한 결과 A는 9,000,000,000원, B는 4,000,000,000원으로 판단된다. 이때 A와 B를 합병하면 토지의 가치는 그림의 C획지와 유사해 질 것으로 판단되며, C의 가치는 18,000,000,000원으로 조사된다. \par
{\small\centering
\begin{tabular}{|c|c|c|c|}
\hline
\makecell{획지} & \makecell{면적(m\textsuperscript{2})} & \makecell{단가(원/m\textsuperscript{2})} & \makecell{가격(원)} \\
\hline
\makecell{A} & \makecell{15,000} & \makecell{600,000} & \makecell{9,000,000,000} \\
\hline
\makecell{B} & \makecell{5,000} & \makecell{800,000} & \makecell{4,000,000,000} \\
\hline
\makecell{C} & \makecell{20,000} & \makecell{900,000} & \makecell{18,000,000,000} \\
\hline
\end{tabular}\par}

\begin{center}
\begin{tikzpicture}[x=0.0125cm,y=0.0125cm,font=\small]
  \draw (0,303) -- (759,303);
  \draw (0,228) -- (759,228);
  \node at (312,263) {도로};
  \draw (75,228) -- (75,2) -- (605,2) -- (605,228);
  \draw (341,228) -- (341,2);
  \draw (190,228) -- (190,116) -- (341,116);
  \draw (215,192) rectangle (319,132);
  \node at (210,210) {B};
  \node at (267,160) {공장1동};
  \node at (472,120) {C};
  \node at (118,43) {A};
\end{tikzpicture}
\end{center}
\textbf{<자료 2> 공장설비에 관한 상황} \par
1. B획지에는 완공된 지 1년이 경과된 공장설비가 소재한다. 이 공장은 전체가 정밀전자부품제조를 위한 특수한 ‘클린룸’ 설비로 구성된 공장이다. \par
2. 재조달원가는 소유자 제시자료에 의거 직접법으로 산정하되, 이의 정산은 기획재정부 계약예규 “예정가격작성기준”에 의거 검토할 예정이며, 부가가치세 및 손해보험료는 고려치 아니하기로 하였다. \par
3. 예정가격작성기준의 적용은 평가대상 설비의 특성상 공사원가계산을 적용하기로 하였고, 참고로 일반관리비율은 5\%, 이윤율은 15\%를 초과할 수 없도록 되어있는 전문 및 기타공사 규정을 적용하여 이것만을 조정하기로 하였다. 한편, 재료비, 노무비, 경비 등은 세부내역을 토대로 검토한 결과 적정한 수준으로 판명되었다. \par
4. 본 공장은 정밀설비로 구성되어 내용년수는 20년으로 조사되며, 지난 1년간 사용상 특이사항이 없어 잔존내용년수는 19년으로 하였고, 잔가율은 10\%로 조사된다. 또한 감가수정은 정밀공장의 고유특성상 정률법을 적용하기로 하였다. \par
5. 지난 1년간 유사 공장(설비)신축에 관한 물가상승률(가격보정지수)은 극히 미미하여 적용치 아니하기로 하였다. \par
\textbf{<자료 3> B획지 소유자 제시 공장의 신축 공사원가계산서(요약)} \par
1. 재료비 : 20억원 \par
2. 노무비 : 10억원 \par
3. 경비 : 10억원 \par
4. 일반관리비 : 4억원 \par
5. 이윤 : 4억원 \par
6. 합계 : 48억원 \par
\end{longquestion}
